% Section 3.3 -- Recovering Signs (role B)
\subsection{Recovering Signs}
\label{sec:B-signs}

After merging, each row $\hat a_j$ equals $s_j c_j\Aone_j$ with an unknown
$c_j>0$ and an unknown sign $s_j\in\{\pm1\}$. The scale $c_j$ is a genuine
symmetry (Proposition~\ref{prop:A-symmetry}), but the sign is not. Using
$\relu(-z)=\relu(z)-z$, replacing $(\Aone_j,\bone_j)$ by
$(-\Aone_j,-\bone_j)$ while keeping $\Atwo_j$ gives
\begin{equation}\label{eq:B-flip}
  \tilde f(x)=f(x)-\Atwo_j\,\pre{j}{x},
\end{equation}
which differs from $f$ by a non-constant affine function whenever
$\Atwo_j\neq0$. Output-layer recovery (Sect.~\ref{sec:A-layer2}) assumes the
first layer is exact, so every sign has to be fixed first.

\paragraph{Pre-image method~\cite[\S4.2]{DBLP:conf/crypto/CarliniJM20}.}
Let $P=\hat A^{+}$ be the pseudo-inverse of the estimated first layer. If
$\done\le\dzero$ and $\hat A$ has full row rank, then $\hat AP=I$, so
moving from $x$ to $x+\tau P_{:,j}$ changes the pre-activation of neuron $j$
by $\tau$ and leaves all others unchanged. For each $j$ we choose a hidden
target $h$ with $h_j=0$ and $|h_k|\in[\tfrac12,1]$ (random signs) for
$k\neq j$, set $x=P(h-\hat b)$, and query $f$ at $x$ and
$x_\pm=x\pm\tau P_{:,j}$ with $\tau=1$. Neuron $j$ is inactive on one side,
where $f$ does not change. If that is the ``$+$'' side, then increasing
$\hat a_j\cdot x$ deactivates the true neuron, so $s_j=-1$ and we negate
$(\hat a_j,\hat b_j)$. Comparing $|f(x_+)-f(x)|$ with $|f(x_-)-f(x)|$ needs no
threshold. The method uses $3\done=30$ queries and fails only if $\Atwo_j=0$,
in which case the sign does not matter.

\paragraph{Least-squares method (ours).}
We also derived a test that uses random queries only. Let
$\hat z=\hat A x+\hat b$ and let $S$ be the set of rows whose sign is
wrong. By $\relu(-z)=\relu(z)-z$,
\begin{equation}\label{eq:B-lstsq}
  f(x)=\sum_{j}w_j\,\relu(\hat z_j)+u\cdot x+u_0,\qquad
  u=-\sum_{j\in S}w_j\hat a_j,\quad u_0=\btwo-\sum_{j\in S}w_j\hat b_j,
\end{equation}
where $w_j=\Atwo_j/c_j$. Equation~\eqref{eq:B-lstsq} holds \emph{exactly for
every} choice of $S$, so we fit $(w,u,u_0)$ by least squares on $n=60$ points
$x\sim\mathcal N(0,4I)$ ($\done+\dzero+1=21$ unknowns). We then solve
$\hat A^{\top}\beta=-u$, which gives $\beta_j=w_j$ for $j\in S$ and $\beta_j=0$
otherwise, and flip row $j$ when $|\beta_j-w_j|<|\beta_j|$. This needs
$\done\le\dzero$, so that $\beta$ is unique, and a full-rank design, which we
check. As a by-product it returns the output layer: $\Atwo=w$ and
$\btwo=u_0+\sum_{j\in S}w_j\hat b_j$.

\paragraph{Comparison.}
Table~\ref{tab:B-signs} starts from the true first layer with random
positive scaling $c_j\sim U[0.1,10]$, random sign flips and relative noise
on every entry, for 100 networks per setting. Without noise and at noise
$10^{-8}$ (larger than the row errors of Sect.~\ref{sec:B-rows}) both
methods recover all signs for all three weight distributions. The
pre-image method needs half the queries, while the least-squares method also
yields an output layer with median error $2^{-42}$--$2^{-41}$ on exact rows. At noise
$10^{-4}$ a few networks fail with both methods. All of them have an
ill-conditioned first layer: the failures have
$\operatorname{cond}(\Aone)$ between $479$ and $1.74\times10^{3}$ (rows
normalised), against medians of $90.7$ for $U[0,1]$ and $30.5$ for
$\mathcal N(0,1)$. A large $\|P\|$ amplifies the row error, so the pre-image
point is no longer close to the hyperplane.

\begin{table}[t]
\centering
\caption{Sign recovery on 100 networks per row. Succ.: networks with all
$10$ signs correct. Last column: median over networks of the output-layer
error obtained as a by-product of our method.}
\label{tab:B-signs}
\setlength{\tabcolsep}{4pt}
\begin{tabular}{llrrrrr}
\toprule
 & & \multicolumn{2}{c}{pre-image~\cite{DBLP:conf/crypto/CarliniJM20}} & \multicolumn{3}{c}{least squares (ours)}\\
\cmidrule(lr){3-4}\cmidrule(lr){5-7}
weights & noise & succ. & queries & succ. & queries & med.\ $\max|f-\hat f|$\\
\midrule
$U[0,1]$ & $0$ & 100/100 & 30 & 100/100 & 60 & $2^{-42}$\\
$U[0,1]$ & $10^{-8}$ & 100/100 & 30 & 100/100 & 60 & $2^{-21}$\\
$U[0,1]$ & $10^{-4}$ & 98/100 & 30 & 98/100 & 60 & $2^{-8}$\\
$U[-1,1]$ & $0$ & 100/100 & 30 & 100/100 & 60 & $2^{-42}$\\
$U[-1,1]$ & $10^{-8}$ & 100/100 & 30 & 100/100 & 60 & $2^{-21}$\\
$U[-1,1]$ & $10^{-4}$ & 100/100 & 30 & 100/100 & 60 & $2^{-7}$\\
$\mathcal N(0,1)$ & $0$ & 100/100 & 30 & 100/100 & 60 & $2^{-41}$\\
$\mathcal N(0,1)$ & $10^{-8}$ & 100/100 & 30 & 100/100 & 60 & $2^{-19}$\\
$\mathcal N(0,1)$ & $10^{-4}$ & 99/100 & 30 & 98/100 & 60 & $2^{-6}$\\
\bottomrule
\end{tabular}

\end{table}
